% jlreq-probe.tex -- LuaLaTeX + jlreq の toolchain 検証 & ground-truth 抽出の土台
%
% ねらい: jlreq クラスは LuaTeX モードで jfm-jlreq.lua を使う = kern が読むのと同じ
%   JFM データ。だから食い違えばエンジン論理 (kern の set-glue vs luatexja) に局在する。
%
% コンパイル (Docker):
%   docker run --rm -v "$PWD":/work -w /work texlive/texlive \
%     lualatex -interaction=nonstopmode jlreq-probe.tex
%
% 見るもの:
%   - jlreq-probe.pdf  : 目視
%   - jlreq-probe.log  : \showbox が吐く hbox の glue/kern ノード (§4.19 の分配)

\documentclass[paper=a5,fontsize=10pt]{jlreq}
\usepackage{luatexja-ruby}

% ボックスの中身をログへ深く吐かせる
\showboxdepth=10
\showboxbreadth=100

\begin{document}

% (1) §4.19 を刺激する行: 約物 (読点・句点・括弧) と和欧混植
日本語の組版、これは「テスト」です。ABC も混植する。

% (2) ルビ: モノルビ (| で区切ると各親字に対応)
\ruby{漢字}{かん|じ}にルビを振る。

% (3) 指定幅の hbox をダンプ = 行長調整の glue 決定を数値で取り出す
%     10 zw に「日本語、テスト。」を詰め込む → 詰め処理が走る
\setbox0=\hbox to 10\zw{日本語、「テスト」。ABC。}
\showbox0

\end{document}
